\section{Main results and proof roadmap}
\label{sec:roadmap}

\paragraph{Proof spine (public theorem chain).}
The public proof is organized around one certification spine and a small number of theorem-level nodes.
\emph{Route~1} is the only load-bearing mass-gap lane:
compact-simple A1 input $\Rightarrow$ one-shot entrance $\Rightarrow$ full fixed-lattice OS gap $\Rightarrow$ continuum vacuum gap.
\emph{Lane~A} is the only load-bearing sharp-local reconstruction lane:
Lane~C flowed state $\Rightarrow$ basis-free exact-dimension quotient package $\Rightarrow$ one-shell / finite-truncation inverse control $\Rightarrow$ finite-cap mixed-correlator bridge $\Rightarrow$ full sharp-local closure with bounded-base cyclicity.
The interface table for Packets~5--7 is Table~\ref{tab:LaneA_packet_interface_map} in Appendix~\ref{app:normalization_crosscheck}, \S\ref{subsec:LaneA_packet_interface_map}: it records that Theorem~\ref{thm:finite_cap_flowed_sector_bridge} is the first mixed-correlator closure node, that Theorem~\ref{thm:LaneA_closure} and Corollary~\ref{cor:LaneA_full_scope} are the capwise-extension and inductive-union nodes, and that Theorem~\ref{thm:LaneA_bounded_base_cyclic} is the first theorem that adds bounded-base cyclicity.
Within Packet~6, the mixed-channel clause reaches its unique downstream theorem-facing closure node at Theorem~\ref{thm:finite_cap_flowed_sector_bridge}(i), read against the compatible family of finite-family packaged bounded capwise states furnished by Corollary~\ref{cor:finite_cap_mixed_family_packaging}.
No universal mixed bounded capwise state is invoked in that step: in Theorem~\ref{thm:finite_cap_flowed_sector_bridge}, the notation $\omega_{\mathrm{cap,bnd}}^{(u_{\mathrm{ren}})}$ is only the local shorthand declared there after the relevant finite family has been fixed.
The natural local citation point for that seam is \hyperref[par:finite_cap_mixed_channel_certification]{the local certification paragraph immediately following Theorem~\ref{thm:finite_cap_flowed_sector_bridge}(i)}, which records the single-slot chain Theorem~\ref{thm:ins_to_OS} $\Rightarrow$ Corollary~\ref{cor:ins_to_OS_bounded_word} $\Rightarrow$ Proposition~\ref{prop:finite_cap_bounded_factor_insertion}, the actual-summand chain Theorem~\ref{thm:finite_truncation_inverse_control} $\Rightarrow$ Lemma~\ref{lem:finite_cap_mixed_channel_reduction} $\Rightarrow$ Theorem~\ref{thm:finite_cap_flowed_sector_bridge}(i), the mixed bounded-family packaging/covariance node Corollary~\ref{cor:finite_cap_mixed_family_packaging}, the fixed-word bookkeeping through $K_W$, $C_W^{\mathrm{ins}}$, and the finite truncation data, and the displayed slotwise-to-functional $o(1)$ handoff in one place.
Appendix~\ref{app:referee_audit}, \S\ref{subsec:packet6_mixed_channel_note}, gives a synchronized theorem-location guide to the same seam together with the same chronology and displayed estimates.
Appendix~\ref{app:referee_audit}, \S\ref{subsec:analytic_gap_packet}, is the short hostile-referee packet for the remaining analytic mass-gap attack line: it isolates the single theorem-level constant route, the finite-cap / positive-flow-time / continuum-limit seam answers, and the faithful-$\rho$ endpoint scope statement in one place.
\emph{Lane~B} is retained only as a downstream dense-core spectral-calculus reformulation and is not used to prove Theorem~\ref{thm:main}.
The fixed-lattice-to-continuum mass-gap transport is frozen locally in \Cref{rem:route1_packet38_transport_certification}: Corollary~\ref{cor:A1_compact_simple_spine} $\Rightarrow$ Proposition~\ref{prop:route1_QE1} $\Rightarrow$ Proposition~\ref{prop:route1_QE2} $\Rightarrow$ Theorem~\ref{thm:Wilson_patch_transport} $\Rightarrow$ Theorem~\ref{thm:continuum_tightness} $\Rightarrow$ Theorem~\ref{thm:continuum_time_OS_upgrade} $\Rightarrow$ Theorem~\ref{thm:phys_gap_noncollapse} $\Rightarrow$ Corollary~\ref{cor:route1_QE3_gap}.
That ledger also records that Theorem~\ref{thm:weak_window_sufficiency} is the only live weak-window certificate on the gap route, that Corollary~\ref{cor:LaneA_bounded_state_restriction} is the first Lane~A theorem-level compatibility input actually consumed there, that this corollary is not the inductive-union theorem, that Theorem~\ref{thm:LaneA_closure}(b) is the finite-cap sharp-local OS-structure input, that Corollary~\ref{cor:LaneA_full_scope} is the separate inductive-union passage, and that Lane~B enters only afterward as a downstream reformulation.
All load-bearing constants needed on these lanes are recorded in the group ledgers of Appendices~\ref{app:A1_compact_simple} and~K.

\paragraph{Certification docket and reading guide.}
Section~\ref{sec:roadmap} together with Appendices~\ref{app:referee_audit} and~\ref{app:executive_proof_sketch} forms one synchronized guide to the public theorem packets.
This section names the public theorem packages, Appendix~\ref{app:referee_audit} begins with Theorem~\ref{thm:proof_spine_certification}, records one primary manuscript location plus one status label for each package, and includes both the Packet~6 seam note \S\ref{subsec:packet6_mixed_channel_note} and the analytic hostile-referee packet \S\ref{subsec:analytic_gap_packet}; Appendix~\ref{app:executive_proof_sketch} follows the same package order without descending into support lemmas.
A new expert can therefore identify the load-bearing theorem stack before reading any detailed proofs.
The front-end Packet~1 certification note, the Packet~3/8 transport ledger, the Packet~5--7 interface map, and the Packet~9/10 provenance ledgers are the synchronized first-pass entry points for the remaining hard modules; Appendix~\ref{app:referee_audit} records the corresponding full-docket release-audit status in theorem-facing form.

\paragraph{Claim-facing packets.}
Packets~1, 9, and~10 are the public claim packets rather than satellite appendices.
Packet~1 places the A1 ultraviolet input on the full stated class of compact connected simple groups and fixes the explicit ultraviolet constants read later by Route~1 for the chosen auxiliary Wilson regularization.
The front-end certification note in Appendix~\ref{app:A1_compact_simple}, \S\ref{subsec:A1_packet1_certification}, records the frozen Packet~1 theorem chain, separates the group-only data from the $(G,\rho)$-indexed ultraviolet constants in Table~\ref{tab:A1_packet1_constant_ledger}, and states explicitly that no separate exceptional-family appendix is load-bearing on the public spine.
Packet~9 promotes the constructive Lane~A / Route~1 output to a local gauge-invariant Yang--Mills theory with OS/Wightman strength, field correspondence, and a non-triviality witness. Appendix~\ref{app:axioms_package} now opens with the provenance ledger in \S\ref{subsec:packet9_provenance_ledger}, which separates the earlier internal Euclidean input, the imported OS reconstruction / analytic-continuation / uniqueness family, and the new theorem-level endpoints packaged there.
Packet~10 fixes the exact theorem-level local-net endpoint of the paper: faithful Wilson regularizations of the same compact simple group are proved to yield canonically isomorphic continuum sharp-local theories, the bounded-region local gauge-invariant net generated by flowed curvature composites is identified as the theorem-level object, and extended-support nonlocal sector data are separated from that local net by theorem rather than by convention. Appendix~\ref{app:global_form_resolution} now opens with the provenance ledger in \S\ref{subsec:packet10_provenance_ledger}, which separates earlier internal universality / descent inputs from the imported OS uniqueness clause used in Theorem~\ref{thm:faithful_Wilson_rep_universality}, the new endpoint theorems, and the AQFT commentary remarks; a short closing remark there then records how that theorem-level endpoint matches the standard local-field reading of the existence problem.

\subsection*{Certified theorem packets (first-pass docket)}
\begin{enumerate}[label=\textbf{Packet \arabic*:},leftmargin=3.35em]
\item Theorem~\ref{thm:A1_compact_simple_internal} and Corollary~\ref{cor:A1_compact_simple_spine}: compact-simple ultraviolet and public group-scope gate; see Appendix~\ref{app:A1_compact_simple}, \S\ref{subsec:A1_packet1_certification}, for the front-end theorem chain and constant ledger.
\item Proposition~\ref{prop:route1_QE1}: one-shot entrance and bounded physical entrance scale.
\item Proposition~\ref{prop:route1_QE2}: tuned-sequence full fixed-lattice OS gap.
\item Theorem~\ref{thm:C3U}, Corollary~\ref{cor:flow_state}, and Corollary~\ref{cor:C3U_Aplus_tuned}: unique flowed continuum state and tuned bounded positive-time base state.
\item Theorem~\ref{thm:LaneA_associated_graded_identity}, Theorem~\ref{thm:LaneA_block_coefficient_theorem}, Theorem~\ref{thm:canonical_block_one_shell_increment}, Theorem~\ref{thm:finite_truncation_inverse_control}, and Corollary~\ref{cor:finite_truncation_SFTE}: basis-free exact-dimension breadth, one-shell transport, and the unconditional finite-truncation inverse / remainder package; see Table~\ref{tab:LaneA_packet_interface_map} for the exact Packet~5--7 role split.
\item Theorem~\ref{thm:finite_cap_flowed_sector_bridge}, Theorem~\ref{thm:LaneA_closure}, and Corollary~\ref{cor:LaneA_full_scope}: first mixed-correlator bridge at finite cap and only then the extension from $\mathcal A_{\mathrm{flow}}$ to the full inductive-union sharp-local algebra.
\item Theorem~\ref{thm:LaneA_bounded_base_cyclic}: bounded-base cyclicity on the reconstructed sharp-local Hilbert space.
\item Theorem~\ref{thm:Wilson_patch_transport}, Theorem~\ref{thm:continuum_tightness}, Theorem~\ref{thm:continuum_time_OS_upgrade}, Theorem~\ref{thm:phys_gap_noncollapse}, and Corollary~\ref{cor:route1_QE3_gap}: Wilson-to-continuum transport and the continuum vacuum gap; see \Cref{rem:route1_packet38_transport_certification} for the local Packet~3/8 ledger naming the only live weak-window certificate, the first Lane~A theorem-level compatibility input, and the later density/core handoff.
\item Theorem~\ref{thm:route1_E2_internal}, Theorem~\ref{thm:R1_HS}, Theorem~\ref{thm:OS_axioms_package}, Theorem~\ref{thm:OS_to_Wightman}, Corollary~\ref{cor:field_correspondence}, and Theorem~\ref{thm:nontriviality_package}: downstream dense-core clustering together with the local-field OS/Wightman existence, explicit field-correspondence, and non-triviality packet; Appendix~\ref{app:axioms_package}, \S\ref{subsec:packet9_provenance_ledger}, is the front-end provenance ledger for that packet.
\item Theorem~\ref{thm:faithful_Wilson_uniform_hypotheses}, Theorem~\ref{thm:faithful_Wilson_rep_universality}, Corollary~\ref{cor:faithful_rho_gap_ledger}, Theorem~\ref{thm:global_form_local_endpoint}, and Corollary~\ref{cor:global_form_main_scope}: faithful-Wilson uniformity for the universality hypotheses, faithful-$\rho$ universality of the continuum local theory, the $\rho$-indexed status of the explicit Route~1 ledger, and the theorem-level local-net endpoint; Appendix~\ref{app:global_form_resolution}, \S\ref{subsec:packet10_provenance_ledger}, is the front-end provenance ledger for that packet.
\end{enumerate}
The only additional certification theorem not on this public docket is \Cref{thm:weak_window_sufficiency}, whose role is not to add new mathematics but to certify that every surviving weak-window use factors through the repaired curvature-sector package of \S\,F.22.2.

\subsection{Main theorem (assembled)}

\begin{theorem}[Local gauge-invariant Yang--Mills QFT and vacuum mass gap for compact connected simple groups]\label{thm:main}
Fix a compact connected simple Lie group $G$, choose any faithful finite-dimensional unitary Wilson representation $\rho:G\to U(d_\rho)$ as an auxiliary ultraviolet regularization, and fix once and for all a weak-window target coupling $u_{\mathrm{ren}}\in(0,u_\ast]$ for the dyadic tuning prescription of Appendix~F.
Appendix~\ref{app:global_form_resolution} proves that any two faithful choices of $\rho$ yield canonically isomorphic continuum sharp-local theories once the same admissible gradient-flow renormalization prescription is imposed, so the theorem-level object depends only on $G$.
We therefore write $\mathfrak A_{\mathrm{loc}}(G)$ for the resulting local gauge-invariant net generated by flowed curvature composites.
The same appendix proves the exact theorem-level local-net endpoint and separates it from any additional nonlocal sector data; the provenance ledger in Appendix~\ref{app:global_form_resolution}, \S\ref{subsec:packet10_provenance_ledger}, records theorem by theorem which inputs are earlier internal universality / descent nodes, where the imported OS uniqueness clause enters, and which remarks are purely AQFT commentary. A short closing remark there then records how that endpoint matches the local-field formulation of Yang--Mills existence.
See Theorem~\ref{thm:faithful_Wilson_uniform_hypotheses}, Theorem~\ref{thm:faithful_Wilson_rep_universality}, Corollary~\ref{cor:faithful_rho_gap_ledger}, Definition~\ref{def:effective_group}, Lemma~\ref{lem:adjoint_local_generators}, Corollary~\ref{cor:vacuum_gap_scope}, Theorem~\ref{thm:global_form_local_endpoint}, and Corollary~\ref{cor:global_form_main_scope}.
The claim-facing appendices are therefore part of the public theorem docket rather than auxiliary resolution notes: Appendix~\ref{app:A1_compact_simple} supplies the compact-simple A1 theorem and traced ultraviolet constants, Appendix~\ref{app:axioms_package} packages the OS/Wightman existence, field-correspondence, and non-triviality theorems and begins with the provenance ledger in \S\ref{subsec:packet9_provenance_ledger}, and Appendix~\ref{app:global_form_resolution} identifies the exact local-net endpoint together with faithful-$\rho$ universality, the uniformity theorem needed by Packet~10, the separation from global-form-sensitive nonlocal sectors, and the provenance ledger in \S\ref{subsec:packet10_provenance_ledger}.
Corollary~\ref{cor:faithful_rho_gap_ledger} records the present quantitative status: the continuum local theory is faithful-$\rho$ universal, but the explicit Route~1 entrance and mass-gap certificate stated here remains indexed by the chosen auxiliary faithful regularization.
Thus, for the chosen auxiliary faithful regularization $\rho$, there exist explicit constants and thresholds recorded in the $(G,\rho)$-ledgers (Appendix~\ref{app:A1_compact_simple}, the audited classical appendices~\ref{app:A1_Sp}--\ref{app:A1_Spin_even}, and Appendix~K) and in the Lane~A/Lane~C ledger
(notably $\beta_\ast(G,\rho)>0$, a one-shot entrance contraction constant $q_\ast(G,\rho)\in(0,1)$, a tuned-sequence physical entrance constant $C_{\mathrm{ent}}(G,\rho)<\infty$,
a trajectory-entry index $N_{\mathrm{ent}}(G,\rho;u_{\mathrm{ren}})\in\mathbb N$, and a block-mixing rate $m_{\mathrm{blk}}(G,\rho):=-\log q_\ast(G,\rho)>0$)
together with the tuned dyadic scaling trajectory associated with $u_{\mathrm{ren}}$ and $\rho$
\[
a_n:=\ell_0\,2^{-n},\qquad \beta_n^{(\rho)}:=\beta(a_n;u_{\mathrm{ren}},G,\rho)\to\infty,\qquad B_n^{(\rho)}:=B(\beta_n^{(\rho)};G,\rho),
\]
such that the following hold:

\begin{enumerate}[label=\textbf{(\Alph*)},leftmargin=2.2em]
\item \textbf{One-shot entrance into the KP/Dobrushin basin (Route~1).}
Let $B(\beta;G,\rho)$ be the explicit one-shot entrance block size from Appendix~F for the chosen faithful ultraviolet regularization $\rho$.
Then for every bare coupling $\beta\ge \beta_\ast(G,\rho)$ the exact one-shot blocked specification at that scale satisfies
\[
\delta\!\bigl(\gamma^{\langle B(\beta;G,\rho)\rangle}\bigr)\le q_\ast(G,\rho)<1,
\]
uniformly in finite volume and admissible boundary conditions.
Along the tuned dyadic trajectory above, the entrance scale is bounded in physical units once the tuned sequence has entered the weak window:
\[
a_n\,B_n^{(\rho)}\le C_{\mathrm{ent}}(G,\rho)\,\ell_0
\qquad (n\ge N_{\mathrm{ent}}(G,\rho;u_{\mathrm{ren}})).
\]
For this fixed faithful regularization $\rho$, the constant $C_{\mathrm{ent}}(G,\rho)$ is uniform in the chosen target $u_{\mathrm{ren}}$; only the entry index $N_{\mathrm{ent}}(G,\rho;u_{\mathrm{ren}})$ depends on that prescription.

\item \textbf{Euclidean exponential clustering on a dense positive-time probe core (downstream Route~1 / Lane~B corollary).}
For the continuum state $\omega$ of \textbf{(C)} and its OS Hilbert space of \textbf{(D)}, there exists a countable positive-time local $*$-algebra
\[
\mathcal G_{\mathrm{ent}}(G)\subset \mathcal A_+
\]
whose centered OS-classes span a dense subspace of the vacuum-orthogonal sector, and for every $F\in\mathcal G_{\mathrm{ent}}(G)$ there is a constant $C_F<\infty$ such that
\[
\big|\omega\big(\Theta \widetilde F\,\tau_s\widetilde F\big)\big|
\le C_F\,e^{-\mathfrak m_\ast(G,\rho;u_{\mathrm{ren}})\, s},
\qquad
\mathfrak m_\ast(G,\rho;u_{\mathrm{ren}}):=\frac{m_{\mathrm{blk}}(G,\rho)}{2\,C_{\mathrm{ent}}(G,\rho)\,\ell_0},
\qquad s\ge 0.
\]

\item \textbf{Continuum local gauge-invariant Yang--Mills QFT with OS/Wightman strength (Lane~A + Appendix~\ref{app:axioms_package}).}
There exists a unique continuum sharp-local state $\omega$ on the full local gauge-invariant sharp-local algebra
\[
\mathfrak A_{\mathrm{loc}}(G)=\varinjlim_{\Delta_{\max}\to\infty}\mathfrak A_{\mathrm{loc}}^{(\le \Delta_{\max})}
\]
generated by flowed curvature composites,
obtained by the $a\downarrow0$ limit followed by the sharp limit $t\downarrow0$ taken \emph{inside} OS pairings of RP-safe representatives.
Appendix~\ref{app:axioms_package} packages the resulting existence theory through Theorem~\ref{thm:OS_axioms_package}, Theorem~\ref{thm:OS_to_Wightman}, Corollary~\ref{cor:field_correspondence}, and Theorem~\ref{thm:nontriviality_package}: the Schwinger functions are Euclidean invariant, symmetric, reflection positive, and tempered; OS reconstruction yields the full sharp-local vacuum representation with $\mathcal A_+$ still cyclic there; standard Osterwalder--Schrader analytic continuation yields a Poincar\'e-covariant Wightman/Haag--Kastler theory whose local fields are the renormalized gauge-invariant local differential polynomials represented by Lane~A; and the canonical centered CP-even scalar witness has nonzero vacuum-created spectral weight.

\item \textbf{Continuum OS/Wightman spectral gap from the full fixed-lattice gap.}
Let $(\mathcal H_{\mathrm{loc}},\Omega_{\mathrm{loc}},H_{\mathrm{loc}})$ be the OS/Wightman reconstruction of Appendix~\ref{app:axioms_package}.
Then
\[
\mathrm{Spec}(H_{\mathrm{loc}})\subset\{0\}\cup[\mathfrak m_\ast(G,\rho;u_{\mathrm{ren}}),\infty),
\qquad
\ker(H_{\mathrm{loc}})=\mathbb C\Omega_{\mathrm{loc}}.
\]
\end{enumerate}

\paragraph{Proof (compressed theorem-level skeleton).}
\begin{enumerate}[label=\textbf{Step \arabic*:},leftmargin=2.6em]
\item \emph{Ultraviolet gate and QE1.}
Corollary~\ref{cor:A1_compact_simple_spine} and Proposition~\ref{prop:route1_QE1} yield \textbf{(A)}.
Packet~1 is therefore both the ultraviolet gate and the public compact-simple group-scope theorem used throughout the paper.
For first-pass auditing, the front of Appendix~\ref{app:A1_compact_simple} records the packet chain, the group-only versus $(G,\rho)$ constant split, and the fact that the exceptional-group case is already closed inside the compact-simple regime theorems.
On this spine, Theorem~\ref{thm:Wilson_patch_transport} is the only return from chart-truncated Lane~C control to the physical Wilson theory, and Theorem~\ref{thm:weak_window_sufficiency} is the only certificate that the repaired curvature-sector weak-window package is the live downstream input from \S\,F.22.2.
\item \emph{Lane~C/Lane~A reconstruction packet.}
Corollary~\ref{cor:flow_state}, Corollary~\ref{cor:C3U_Aplus_tuned}, Theorem~\ref{thm:LaneA_associated_graded_identity}, Theorem~\ref{thm:LaneA_block_coefficient_theorem}, Theorem~\ref{thm:finite_truncation_inverse_control}, Theorem~\ref{thm:finite_cap_flowed_sector_bridge}, Theorem~\ref{thm:LaneA_closure}, Corollary~\ref{cor:LaneA_bounded_state_restriction}, Corollary~\ref{cor:LaneA_full_scope}, and Theorem~\ref{thm:LaneA_bounded_base_cyclic} supply the flowed continuum state, solve the exact-dimension breadth problem, prove the finite-cap mixed-channel bridge, extend the Lane~C state capwise to the sharp-local algebra, record the bounded-state compatibility later consumed on the QE3 seam, pass separately to the full inductive union, and preserve cyclicity of the bounded positive-time base algebra in the reconstructed sharp-local Hilbert space.
At finite cap, Theorem~\ref{thm:finite_cap_flowed_sector_bridge}(i) is the unique downstream theorem-facing closure node for the mixed-channel clause, read against the compatible family of finite-family packaged bounded capwise states furnished by Corollary~\ref{cor:finite_cap_mixed_family_packaging}.
The natural local citation point for that seam is \hyperref[par:finite_cap_mixed_channel_certification]{the local certification paragraph immediately following Theorem~\ref{thm:finite_cap_flowed_sector_bridge}(i)}; it records the same two proof chains, the mixed bounded-family packaging/covariance node Corollary~\ref{cor:finite_cap_mixed_family_packaging}, the fixed-word bookkeeping through $K_W$, $C_W^{\mathrm{ins}}$, and the finite truncation data, and the displayed slotwise-to-functional $o(1)$ handoff that Appendix~\ref{app:referee_audit}, \S\ref{subsec:packet6_mixed_channel_note}, summarizes in theorem-location form.
This is the only sharp-local reconstruction packet later used by QE3 and Appendix~\ref{app:axioms_package}.
\item \emph{QE2 and QE3.}
Proposition~\ref{prop:route1_QE2}, Theorem~\ref{thm:Wilson_patch_transport}, Theorem~\ref{thm:continuum_tightness}, Theorem~\ref{thm:continuum_time_OS_upgrade}, Theorem~\ref{thm:phys_gap_noncollapse}, and Corollary~\ref{cor:route1_QE3_gap} convert the tuned full fixed-lattice OS gap into the continuum sharp-local vacuum gap on the full reconstructed Hilbert space.
The frozen local certification point for that fixed-lattice-to-continuum step is \Cref{rem:route1_packet38_transport_certification}, which also records Theorem~\ref{thm:weak_window_sufficiency} as the only live weak-window certificate, Corollary~\ref{cor:LaneA_bounded_state_restriction} as the first Lane~A theorem-level compatibility input on the transport route, Theorem~\ref{thm:LaneA_closure}(b) as the finite-cap sharp-local OS-structure input, Corollary~\ref{cor:LaneA_full_scope} as the separate inductive-union passage, and Corollary~\ref{cor:sharp_base_dense_in_full_H} together with Lemma~\ref{lem:QE3_core_topology} as the later density/core handoff.
This yields \textbf{(D)}.
The dense-core Euclidean-clustering / Lane~B reformulation in \textbf{(B)} comes only afterward through Theorem~\ref{thm:route1_laneB_package}.
\item \emph{Claim packets: axioms, non-triviality, and the local-net endpoint.}
Theorem~\ref{thm:OS_axioms_package}, Theorem~\ref{thm:OS_to_Wightman}, Corollary~\ref{cor:field_correspondence}, Theorem~\ref{thm:nontriviality_package}, Theorem~\ref{thm:faithful_Wilson_uniform_hypotheses}, Theorem~\ref{thm:faithful_Wilson_rep_universality}, Corollary~\ref{cor:faithful_rho_gap_ledger}, Theorem~\ref{thm:global_form_local_endpoint}, and Corollary~\ref{cor:global_form_main_scope} package the Euclidean existence theorem, the standard Wightman/Haag--Kastler reconstruction, the explicit field-correspondence statement, the non-triviality witness, the faithful-Wilson uniform-hypothesis theorem, the faithful-Wilson universality theorem, the present quantitative-ledger status, and the exact local-net endpoint. Appendices~\ref{app:axioms_package} and~\ref{app:global_form_resolution} now front-load the corresponding provenance ledgers in \S\ref{subsec:packet9_provenance_ledger} and \S\ref{subsec:packet10_provenance_ledger}.
Together these claim-facing packets give \textbf{(C)} and fix the theorem-level local-net endpoint of the result.
\end{enumerate}

\begin{remark}[Frozen lane roles and exact local-net endpoint]\label{rem:frozen_spine_scope}
For a first-pass audit of \Cref{thm:main}, it is enough to keep the ten public packets listed above together with the certification theorem \Cref{thm:weak_window_sufficiency} in view.
The proof of \Cref{thm:main}(D) runs only through the primary Route~1 spine named in Steps~1--3.
Lane~A is load-bearing exactly where Appendix~F needs the extension from $\mathcal A_{\mathrm{flow}}$ to the full local gauge-invariant sharp-local algebra, the continuum-time OS semigroup, bounded-base cyclicity on the reconstructed Hilbert space, and the separate temperedness verification for the full inductive-union Schwinger family.
Table~\ref{tab:LaneA_packet_interface_map} keeps the theorem roles literal there: the finite-cap bridge is the first mixed-correlator closure node, the inductive-union passage occurs only at Corollary~\ref{cor:LaneA_full_scope}, and bounded-base cyclicity appears only at Theorem~\ref{thm:LaneA_bounded_base_cyclic}; the later Packet~9-facing Lane~A exports are then Theorem~\ref{thm:LaneA_temperedness_verification} and Theorem~\ref{thm:LaneA_nontriviality_witness}.
Lane~B is not part of that proof: it is retained only as a downstream reformulation of the already established vacuum-gap statement from \Cref{thm:main}(B).
The term ``mass gap'' throughout the paper therefore refers to the vacuum Hamiltonian gap in the OS reconstruction of the local gauge-invariant sharp-local net, not to additional superselection sectors created by nonlocal defects; Appendix~\ref{app:global_form_resolution} fixes that theorem-level scope precisely. Appendices~\ref{app:A1_compact_simple}, \ref{app:axioms_package}, and~\ref{app:global_form_resolution} are accordingly load-bearing claim packets, not downstream commentary.
\end{remark}

\paragraph{Route~1 proof hygiene.}
The QE1 entrance theorem is a finite-depth ultraviolet statement: it uses only finite-volume local analytic input together with exact Wilson recursion/transport closure and does \emph{not} invoke any infrared clustering, infinite-volume limit, or spectral information.
The repaired weak-window package enters only through the tuned physical-scale normalization encoded in $N_{\mathrm{ent}}(G,\rho;u_{\mathrm{ren}})$ and $C_{\mathrm{ent}}(G,\rho)$; the exact downstream dependence is certified separately by \Cref{thm:weak_window_sufficiency}.
\end{theorem}

\begin{remark}[Load-bearing non-triviality witness]\label{rem:nontriviality_witness}
Appendix~\ref{app:axioms_package} proves non-triviality at theorem level rather than by appeal to the sign of the renormalized coupling.
Theorem~\ref{thm:LaneA_nontriviality_witness} fixes a canonical centered CP-even scalar field with unit OS norm on a reference positive-time test function,
and Theorem~\ref{thm:nontriviality_package} upgrades that witness to a nonzero Euclidean/Wightman two-point function and a nonzero K\"all\'en--Lehmann spectral measure supported above the vacuum gap.
Accordingly the local gauge-invariant sharp-local theory constructed here is not the zero theory and not the identity theory in its vacuum sector.
\end{remark}

\subsection{Global form: \texorpdfstring{$\Spin$}{Spin} versus \texorpdfstring{$\SO$}{SO} and the meaning of the gap}
\label{subsec:spin_vs_so}

The two simply-connected series $\Spin(2n{+}1)$ and $\Spin(2n)$ admit nontrivial central quotients.
The statements collected below are the Section~\ref{sec:roadmap} front end of Appendix~\ref{app:global_form_resolution}: they are included here so that the title, abstract, and main theorem already name the exact theorem-level endpoint.
Since our lattice input is a Wilson plaquette weight built from the \emph{defining} representation $\rho$
(\Cref{prop:A1_template}), it is important to pin down the exact global form of the gauge group
that is actually seen by the construction.

\begin{definition}[Effective gauge group determined by the plaquette action]
\label{def:effective_group}
Let $G$ be a compact Lie group and let $\rho:G\to U(d_\rho)$ be the representation used to define the Wilson plaquette weight.
Set
\begin{equation}
\label{eq:kernel_K}
K_\rho\ :=\ \ker(\rho)\ \subset\ Z(G),
\qquad
G_{\mathrm{eff}}\ :=\ G/K_\rho.
\end{equation}
\end{definition}

\begin{lemma}[Center-blindness and descent to the quotient]
\label{lem:center_blindness_descent}
With the Wilson action defined from $\rho$, the one-plaquette weight, the full Wilson lattice gauge measure,
and every gauge-invariant observable generated by flowed local composites of the curvature (hence in particular
every observable in the sharp local algebra $\mathcal A_{\mathrm{loc}}$) depend on $G$ only through the quotient $G_{\mathrm{eff}}$.
Equivalently, the resulting Schwinger functions and the reconstructed sharp local OS/Wightman theory are canonically
identified for $(G,\rho)$ and $(G_{\mathrm{eff}},\rho_{\mathrm{desc}})$.
\end{lemma}

\begin{proof}
The plaquette action density has the form
$A_\rho(U)=1-\frac{1}{d_\rho}\Re\Tr_\rho(U)$.
If $k\in K_\rho$ then $\rho(kU)=\rho(U)$, hence $A_\rho(kU)=A_\rho(U)$.
Therefore the Wilson weight is constant on cosets of $K_\rho$ and the lattice measure is invariant under pointwise multiplication
of link variables by elements of $K_\rho$.
Any gauge-invariant observable built from $\rho$ (and any flowed local composite built from the curvature) is likewise $K_\rho$-invariant,
so its expectation depends only on the induced variables in $G/K_\rho$.
Passing to the continuum limit and OS reconstruction preserves these identities.
\end{proof}

\begin{lemma}[Adjoint-local generators of the sharp local algebra and insensitivity to global form]
\label{lem:adjoint_local_generators}
The sharp local algebra $\mathcal A_{\mathrm{loc}}$ constructed in \S\ref{sec:sharp_local} is generated by renormalized local composite
insertions obtained as small-flow-time limits of \emph{flowed curvature composites} (elements of $\mathcal A_{\mathrm{flow}}$),
i.e.\ smeared fields of the form
\[
\int d^4x\ f(x)\,P\!\big(\{(D^{k}G_{\mu\nu})(t,x)\}_{|k|\le k_0,\mu,\nu}\big),
\qquad t>0,
\]
where $G_{\mu\nu}(t,x)\in\mathfrak g$ is the flowed curvature and $P$ is an $\Ad$-invariant polynomial (equivalently, a linear combination of
complete contractions using $\Ad$-invariant tensors; in particular traces in finite-dimensional representations).
In particular, every generator is \emph{adjoint-local}: it depends on the gauge field only through the associated $\Ad$-bundle,
so it is invariant under pointwise multiplication by central gauge transformations.
Consequently, for any central subgroup $K\le Z(G)$, the local net generated by $\mathcal A_{\mathrm{loc}}$ is canonically identified for $G$ and $G/K$,
and $\mathcal A_{\mathrm{loc}}$ cannot distinguish different global forms with the same Lie algebra.
\end{lemma}

\begin{proof}
By definition of the flowed probe algebra $\mathcal A_{\mathrm{flow}}$ in Lane~C (see \S\ref{sec:rg}, in particular the definition preceding
Lemma~\ref{lem:F_to_J}), its generators are gauge-invariant flowed curvature composites built from $G_{\mu\nu}(t,x)$ and covariant derivatives.
Lane~A identifies each renormalized sharp local insertion as a $t\downarrow 0$ limit of inverse-flow linear combinations of such probes
(Theorem~\ref{thm:LaneA_closure}), hence $\mathcal A_{\mathrm{loc}}$ is generated by these curvature composites.

Under a (time-slice) gauge transformation $g(x)\in G$, the flowed curvature transforms as
$G_{\mu\nu}(t,x)\mapsto \Ad(g(x))\,G_{\mu\nu}(t,x)$, and likewise for covariant derivatives.
Since $P$ is $\Ad$-invariant, the composite is gauge-invariant.
If $k(x)\in K\subset Z(G)$ is a central gauge transformation, then $\Ad(k(x))=\mathrm{id}_{\mathfrak g}$, so every such generator is unchanged.
Because the adjoint action is identical for $G$ and $G/K$, the collection of adjoint-local composites and the $\ast$-algebra they generate
are canonically identified for the two global forms.
\end{proof}

\begin{corollary}[Scope of the vacuum gap statement]
\label{cor:vacuum_gap_scope}
The vacuum spectral gap asserted in \Cref{thm:main}(D) is a statement about the OS vacuum representation
of the local gauge-invariant net generated by $\mathcal A_{\mathrm{loc}}$ (\Cref{lem:adjoint_local_generators}).
In particular, it does \emph{not} assert a gap in other representations or superselection sectors that may be generated
by nonlocal line/surface operators (e.g. Wilson lines in spinor representations, 't~Hooft operators, or other defect insertions).
Appendix~\ref{app:global_form_resolution} shows that this local-net interpretation is the theorem-level endpoint used in the paper.
\end{corollary}


\begin{remark}[What this means for the two \texorpdfstring{$\Spin$}{Spin} families]
\label{rem:spin_vs_so_meaning}
For $G=\Spin(2n{+}1)$ and $G=\Spin(2n)$, the defining (vector) representation factors through the covering map
$\Spin(\cdot)\twoheadrightarrow \SO(\cdot)$, hence $K_\rho=\{\pm\mathbf 1\}$ and $G_{\mathrm{eff}}\cong\SO(2n{+}1)$ or $\SO(2n)$.
Accordingly, the sharp local net constructed in \Cref{thm:main} is the \emph{local} Yang--Mills theory associated with the common Lie algebra
and is insensitive to the global form (simply-connected vs adjoint) at the level of strictly local gauge-invariant observables (\Cref{lem:adjoint_local_generators}).

The difference between $\Spin$ and $\SO$ is encoded in \emph{nonlocal} line/surface operators (e.g. Wilson lines in spinor representations and
their associated superselection structure).
This manuscript constructs the OS/Wightman theory generated by the sharp local algebra $\mathcal A_{\mathrm{loc}}$ (\Cref{lem:adjoint_local_generators});
the mass gap assertion in \Cref{thm:main}(D) refers to the spectral gap of the Hamiltonian in the vacuum representation obtained from OS reconstruction of this local net.
Additional representations of the same local net (if any) correspond to distinct superselection sectors and are not asserted to be covered by the present gap statement.
\end{remark}

\subsection{One-page dependency graph}
\label{subsec:dependency_graph}

The proof is organized in two coupled lanes.

\begin{itemize}
\item \textbf{Lane A (construction):}
reflection positivity together with exact Wilson recursion and Wilson-to-patched transport (\S\ref{sec:reflection});
continuum flow-bundle state by dyadic gradient-flow RG (\S\ref{sec:rg});
sharp local operator algebra and OS axioms (\S\ref{sec:sharp_local}).

\item \textbf{Lane B (gap):}
Route~1 mass-gap theorem along the fixed-flow-time scaling trajectory (Appendix~\ref{app:route1});
OS intake and transfer of exponential clustering to a spectral gap in the sharp local limit (\S\ref{sec:mass_gap_intake}).
\end{itemize}

All cross-lane interactions occur through a finite list of explicit constants recorded in \S\ref{sec:roadmap_constants}.

\subsection{Explicit constant status}
\label{sec:roadmap_constants}

For auditability, every parameter constraint is stated as an inequality between named constants.
The load-bearing family-dependent constants are:
\begin{equation}
\label{eq:constant_tuple}
(b_0(G),\ \kappa_G,\ \lambda^*_G,\ \rho_G,\ K_{0,G},\ c_{\mathrm{A1}}(G,\rho),\ \beta_{\mathrm{A1}}(G,\rho)).
\end{equation}
Here $c_{\mathrm{A1}}(G,\rho),\beta_{\mathrm{A1}}(G,\rho)$ are the A1 one-plaquette crossover constants for the chosen auxiliary Wilson regularization (\Cref{prop:A1_template}).

\begingroup
\small
\begin{longtable}{@{}
>{\RaggedRight\arraybackslash}p{0.18\textwidth}
>{\RaggedRight\arraybackslash}p{0.23\textwidth}
>{\RaggedRight\arraybackslash}p{0.18\textwidth}
>{\RaggedRight\arraybackslash}p{0.33\textwidth}
@{}}
\toprule
\textbf{Input} & \textbf{Role} & \textbf{Where proved} & \textbf{Notes} \\
\midrule
\endhead
$b_0(G)>0$ & step-scaling sign & Appendix~F & one-loop coefficient in the dyadic GF scheme \\
$\kappa_G$ & one-plaquette coercivity & Appendix~F & curvature localization/BL inequality \\
$\lambda^*_G$ & tree-gauge Hessian gap & Appendix~E & Sturm--Habicht root isolation; universal combinatorial constant \\
$\rho_G<1$ & RG contraction & \S\ref{sec:rg} & seminorm contraction factor \\
$K_{0,G}$ & heat-kernel constants & Appendices B--D & only mild family dependence \\
$c_{\mathrm{A1}}(G,\rho)$, $\beta_{\mathrm{A1}}(G,\rho)$ & one-plaquette crossover (A1) & Appendix~\ref{app:A1_compact_simple} & compact-simple harmonic analysis for the chosen auxiliary Wilson regularization; Appendices~G--I remain audited classical examples \\
\bottomrule
\end{longtable}
\endgroup

\begin{remark}[Why the A1 shape is unavoidable]
\label{rem:no_go_heatkernel_generic}
A uniform ``heat-kernel'' bound of the form
$|a_{\pi}(\beta,\alpha)|\le \exp(-c\,C_2(\pi)/\beta)$ for all irreps $\pi$ and fixed $\beta$
fails already in rank-one nonabelian settings: along an embedded root subgroup the coefficient decay necessarily crosses over from
Gaussian-in-$C_2/\beta$ behavior (small representations) to Gaussian-in-$\sqrt{C_2}$ behavior (large representations).
Since every compact connected simple Lie group contains rank-one root subgroups,
a uniform bound of purely Gaussian-in-$C_2/\beta$ type is not available at arbitrarily large highest weights.
The crossover denominator $\beta+\sqrt{C_2(\pi)}$ in \eqref{eq:A1_canonical} is the sharp substitute compatible with both regimes.
\end{remark}

